
\documentclass[12pt,a4paper]{report}

\usepackage[utf8]{inputenc}
\usepackage{geometry}
\usepackage{graphicx}
\usepackage{booktabs}
\usepackage{array}
\usepackage{float}
\usepackage{hyperref}
\usepackage{setspace}

\geometry{margin=1in}
\onehalfspacing

\title{
    \textbf{ONLINE RETAIL SALES DATA PIPELINE}\\
    \large Mini Project Report
}
\author{Manmitha}
\date{November 2026}

\begin{document}

\maketitle

\tableofcontents
\newpage

\chapter{Introduction}

\section{Project Overview}

The Online Retail Sales Data Pipeline is a reusable
ETL (Extract, Transform, Load) project developed to
clean and process online retail transaction data.

The project uses Python and Pandas to extract data
from an Excel file, perform data cleaning, generate
business-related features, validate the processed
dataset, and export the results in CSV and Parquet
formats.

\section{Project Objective}

The main objectives are:

\begin{itemize}
    \item Load and process online retail transaction data.
    \item Remove invalid and duplicate records.
    \item Handle missing values appropriately.
    \item Generate useful customer and sales features.
    \item Develop reusable ETL modules.
    \item Implement logging and exception handling.
    \item Use checkpointing to reuse processed data.
    \item Validate the final dataset.
    \item Track dataset versions using DVC.
    \item Export data in CSV and Parquet formats.
\end{itemize}

\chapter{Dataset Description}

\section{Dataset Information}

The project uses the Online Retail dataset from
the UCI Machine Learning Repository.

\begin{table}[H]
\centering
\begin{tabular}{ll}
\toprule
\textbf{Attribute} & \textbf{Value}\\
\midrule
Dataset Name & Online Retail\\
Original File & Online Retail.xlsx\\
Original Records & 541,909\\
Original Columns & 8\\
Final Records & 526,052\\
Final Columns & 17\\
Output Formats & CSV and Parquet\\
\bottomrule
\end{tabular}
\caption{Dataset Information}
\end{table}

\section{Original Dataset Attributes}

The original dataset contains the following columns:

\begin{itemize}
    \item InvoiceNo: Invoice number.
    \item StockCode: Product identification code.
    \item Description: Product description.
    \item Quantity: Number of products purchased.
    \item InvoiceDate: Date and time of transaction.
    \item UnitPrice: Price per unit.
    \item CustomerID: Customer identifier.
    \item Country: Customer's country.
\end{itemize}

\chapter{Methodology}

\section{ETL Pipeline}

The project follows three main stages:

\begin{enumerate}
    \item \textbf{Extract:} Load transaction records
    from the Excel dataset.
    \item \textbf{Transform:} Clean the data, handle
    invalid records, and generate additional features.
    \item \textbf{Load:} Save the processed dataset
    in CSV and Parquet formats.
\end{enumerate}

\section{Data Cleaning}

The following operations were performed:

\begin{itemize}
    \item Removed cancelled invoices.
    \item Excluded transactions with negative quantities.
    \item Removed records with negative unit prices.
    \item Removed duplicate records.
    \item Replaced missing product descriptions with
    \texttt{Unknown}.
    \item Retained missing CustomerID records and
    created an indicator column.
    \item Retained zero-price transactions and
    created an indicator column.
    \item Converted InvoiceDate into datetime format.
\end{itemize}

\section{Data Cleaning Summary}

\begin{table}[H]
\centering
\begin{tabular}{lr}
\toprule
\textbf{Operation} & \textbf{Records}\\
\midrule
Original records & 541,909\\
Cancelled invoices removed & 9,288\\
Negative quantity records excluded & 1,336\\
Negative unit-price records removed & 2\\
Duplicate records removed & 5,231\\
\midrule
Final processed records & 526,052\\
\bottomrule
\end{tabular}
\caption{Data Cleaning Summary}
\end{table}

\section{Feature Engineering}

The following features were generated:

\begin{itemize}
    \item TotalAmount
    \item InvoiceMonth
    \item InvoiceYear
    \item InvoiceDay
    \item CustomerTotalOrders
    \item AverageOrderValue
    \item TotalCustomerSpending
\end{itemize}

The transaction amount is calculated as:

\[
\text{TotalAmount} =
\text{Quantity} \times \text{UnitPrice}
\]

Customer-level features are calculated using
available CustomerID values.

\chapter{Implementation}

\section{Technology Stack}

The project uses:

\begin{itemize}
    \item Python
    \item Pandas
    \item Python Calamine
    \item PyArrow
    \item Jupyter Notebook
    \item Git and GitHub
    \item DVC
\end{itemize}

\section{Reusable Modules}

The source code is organized into reusable modules:

\begin{itemize}
    \item \texttt{data\_loader.py}: Loads Excel data.
    \item \texttt{data\_cleaning.py}: Performs data cleaning.
    \item \texttt{feature\_engineering.py}: Generates
    additional features.
    \item \texttt{validation.py}: Validates processed data.
    \item \texttt{pipeline.py}: Coordinates the ETL process.
\end{itemize}

\section{Logging and Exception Handling}

Logging is implemented to record pipeline execution
and processing information.

Exception handling is used to manage errors during
data processing and pipeline execution.

\section{Checkpointing}

Checkpointing allows the pipeline to reuse previously
processed intermediate data when available, reducing
unnecessary repeated processing.

\section{Version Tracking}

Git is used for source code version tracking, while
DVC is used to track dataset versions.

The processed dataset is tracked using DVC.

\chapter{Results and Validation}

\section{Final Dataset}

The final processed dataset contains 526,052 records
and 17 columns.

The dataset is exported in two formats:

\begin{itemize}
    \item CSV
    \item Parquet
\end{itemize}

\section{Data Validation Results}

The following validation checks were performed:

\begin{table}[H]
\centering
\begin{tabular}{ll}
\toprule
\textbf{Validation Check} & \textbf{Result}\\
\midrule
Dataset shape & 526,052 $\times$ 17\\
Missing required columns & 0\\
Duplicate rows & 0\\
Negative quantities & 0\\
Negative unit prices & 0\\
Cancelled invoices & 0\\
Invalid invoice date type & False\\
TotalAmount calculation errors & 0\\
\bottomrule
\end{tabular}
\caption{Final Data Validation Results}
\end{table}

\section{CSV and Parquet Comparison}

\begin{table}[H]
\centering
\begin{tabular}{lr}
\toprule
\textbf{Format} & \textbf{File Size}\\
\midrule
CSV & 69.00 MB\\
Parquet & 4.25 MB\\
\bottomrule
\end{tabular}
\caption{Storage Comparison}
\end{table}

Parquet requires approximately 93.84\% less storage
than the CSV output in this project.

\chapter{Conclusion}

The Online Retail Sales Data Pipeline was successfully
developed to process and validate online retail
transaction data.

The project implements data extraction, cleaning,
feature engineering, validation, logging,
exception handling, checkpointing, and version tracking.

The final dataset contains 526,052 records and
17 columns and is available in CSV and Parquet formats.

This project demonstrates the development of a
reusable data processing pipeline and provides
a foundation for further retail sales and
customer analysis.

\end{document}
